\section{Ordered states in the residue carrier}
\label{sec:ordered-residue-coefficients}

A residue acts on a differential form, while an ordered bar face multiplies adjacent states.
These operations can be combined in two ways.
One can carry an entire bar complex unchanged along every normal stratum, or attach a multiplication map to the stratum reached by a collision.
The first construction inherits the unfiltered overlap homotopies.
The second requires the coefficient diagram itself; its behavior is already different for an algebra with square-zero augmentation ideal.
We construct both complexes and their coordinate maps.

All coefficients in this section are over $\mathbb C$.
Tensor products of vector spaces are algebraic.
For sheaves, tensor product means the sheaf associated to the algebraic tensor product on local sections.
Thus each germ is a finite sum of tensors; no completion is used.
Normal charts and their base projection satisfy the hypotheses of Section~\ref{sec:cech-map}.

\subsection{A fixed ordered bar complex as coefficients}
\label{subsec:fixed-bar-coefficients}

Let $H$ be an augmented differential graded associative algebra and $I$ its augmentation ideal.
Use the desuspension $t$ of degree $-1$ and set
\[
Q=\bigoplus_{n\geq0}(tI)^{\otimes n},\qquad b=b_1+b_\mu.
\]
The empty word is a copy of $\mathbb C$.
For homogeneous $a_i\in I$, the word $[a_1|\cdots|a_n]$ has degree
$\sum_i|a_i|-n$.
The formulas~\eqref{eq:bar-corestrictions}--\eqref{eq:bar-binary-sign} define $b$;
Proposition~\ref{prop:bar-determinant-comparison} proves $b^2=0$ by the derivation and associativity identities.
In particular, this definition includes the changing word length under an adjacent bar face.

Let $\mathcal A=A_E(\Omega_Z^\bullet)$ be the sheaf of normal complexes from~\eqref{eq:normal-total} and let
$\mathcal B=p_*\mathcal R_Y$.
View $Q$ as a constant sheaf of complexes on $Z$ and put
\begin{equation}
\label{eq:fixed-bar-residue-carriers}
\mathcal A_Q=Q\otimes\mathcal A,\qquad
\mathcal B_Q=Q\otimes\mathcal B.
\end{equation}
The differential on either tensor product is
\begin{equation}
\label{eq:fixed-bar-total-differential}
D_Q(w\otimes\xi)=bw\otimes\xi+(-1)^{|w|}w\otimes D\xi,
\end{equation}
using $D_A$ on $\mathcal A$ and $D$ on $\mathcal B$.
Thus the coefficient differential includes the bar multiplications, whereas a residue changes the normal stratum and leaves $w$ unchanged.
The word length and the subset of normal divisors are independent indices in~\eqref{eq:fixed-bar-residue-carriers}.

\begin{lemma}[Tensor signs for comparison maps]
\label{lem:tensor-residue-homotopies}
Let $f:A\to B$ have degree $a$ and write
$\partial f=d_Bf-(-1)^a f d_A$.
Its tensor extension is
\[
f^Q(w\otimes\xi)=(-1)^{a|w|}w\otimes f\xi.
\]
Then $\partial(f^Q)=(\partial f)^Q$, and $(gf)^Q=g^Qf^Q$ whenever the composition is defined.
\end{lemma}

\begin{proof}
Write $q=|w|$.
The terms containing $bw$ in $\partial(f^Q)$ have coefficients
$(-1)^{aq}$ and $-(-1)^{a+a(q+1)}$, whose sum is zero.
The remaining terms are
\[
(-1)^{(a+1)q}w\otimes\bigl(d_Bf-(-1)^af d_A\bigr)\xi.
\]
This is $(\partial f)^Q$.
For composition, the two signs multiply to
$(-1)^{(|f|+|g|)q}$, as required.
\end{proof}

\begin{proposition}[Coherent descent with ordered bar coefficients]
\label{prop:bar-residue-descent}
The differentials~\eqref{eq:fixed-bar-total-differential} square to zero.
The full carrier $\mathcal B_Q$ has strict coordinate pullbacks
$\mathrm{id}_Q\otimes g^*$.
The normal embeddings and determinant transports extend to chain maps
$j_i^Q=\mathrm{id}_Q\otimes j_i$ and $P_{ij}^Q=\mathrm{id}_Q\otimes P_{ij}$.
For $m>0$, their overlap homotopies are
\begin{equation}
\label{eq:bar-overlap-homotopy}
H_{ij}^Q(w\otimes\xi)=(-1)^{|w|}w\otimes H_{ij}\xi.
\end{equation}
They satisfy
\[
D_QH_{ij}^Q+H_{ij}^QD_Q=j_j^Q-j_i^QP_{ij}^Q,
\qquad
H_{ik}^Q=H_{jk}^Q+H_{ij}^QP_{jk}^Q.
\]
The source $\mathcal A_Q$ is contractible by $h^Q(w\otimes\xi)=(-1)^{|w|}w\otimes h\xi$.
\end{proposition}

\begin{proof}
In the square of~\eqref{eq:fixed-bar-total-differential}, the terms involving both $b$ and $D$
have coefficients $(-1)^{|w|}$ and $(-1)^{|w|+1}$.
They cancel; the remaining terms vanish by $b^2=D^2=0$.
The strict pullbacks are chain maps by Proposition~\ref{prop:strict-residue-descent} and their identity action on $Q$.
Their composition follows from composition of pullbacks.
Apply Lemma~\ref{lem:tensor-residue-homotopies} to $j_i$, $P_{ij}$, $H_{ij}$ and $h$.
The identities of Theorem~\ref{thm:strict-homotopy-cocycle} and Lemma~\ref{lem:normal-contractible}
then give every displayed identity, including the contraction.
\end{proof}

On a finite base cover, the actual total comparison sends a homogeneous element
$w\otimes\xi$ of \v Cech degree $p$ to the two components
\[
j_{i_0}^Q(w\otimes\xi),\qquad
(-1)^{p+|w|}w\otimes H_{i_0i_1}\xi.
\]
The second expression uses the tail indices of~\eqref{eq:cech-comparison}.
Its sign includes both the \v Cech degree and the degree of the ordered word.
The proof of Proposition~\ref{prop:cech-comparison}, applied to the displayed overlap identities,
proves the total chain-map equation.

The filtered logarithmic homotopies extend by the same rule:
$K^Q(w\otimes\xi)=(-1)^{|w|}w\otimes K\xi$ and
$L^Q(w\otimes\xi)=w\otimes L\xi$.
Thus, for two normal directions and additive logarithms, the three total components have coefficients
$1$, $(-1)^{p+|w|}$, and $-1$ multiplying $j$, $K$, and $L$, respectively.
This supplies the filtered total map, including its secondary homotopy.

For example, take $H=\mathbb C[\epsilon]/(\epsilon^2)$ in degree zero, with zero differential.
The closed bar word $[\epsilon]$ has degree $-1$.
With one normal coordinate changed by $y=e^z x$, the logarithmic homotopy sends
\[
[\epsilon]\otimes e\longmapsto-[\epsilon]\otimes z.
\]
The form part of $D_Q$ also has a minus sign on this word.
It therefore sends the right side to $[\epsilon]\otimes dz$, the difference of the two normal embeddings.
Omitting the sign in~\eqref{eq:bar-overlap-homotopy} would reverse this equation.

Neither tensoring with $Q$ nor taking its bar differential removes the scalar limitations.
The inclusion of the empty word and the projection onto it are cochain maps
$\mathbb C\to Q\to\mathbb C$ whose composite is the identity.
Tensoring these maps with every stratum gives a filtered retraction onto the scalar complexes.
Any filtered overlap or secondary homotopy for the specified tensor comparisons would therefore induce the corresponding scalar one.
The $\mathbb C^*$ period obstruction of Example~\ref{ex:period-obstruction} and the nonzero constants of
Theorem~\ref{thm:filtered-triple-criterion} still obstruct those homotopies.
The same retraction carries the nonzero constant class in $H^0(\mathcal B)$ into $H^0(\mathcal B_Q)$;
its existence rules out a quasi-isomorphism from the contractible source $\mathcal A_Q$ when $m>0$.

\subsection{Multiplication attached to a collision stratum}
\label{subsec:state-changing-residues}

We now let a residue multiply the states at its two adjacent blocks.
Fix an ordered string of $n\geq1$ positions and its $m=n-1$ cuts
$E=\{1,\ldots,n-1\}$.
For $S\subseteq E$, declare positions $i$ and $i+1$ equivalent when $i\in S$.
The equivalence classes are consecutive intervals, ordered from left to right.
Call them the blocks of $S$ and write $b(S)=n-|S|$ for their number.
Assign one state in $I$ to each block:
\begin{equation}
\label{eq:interval-coefficient-system}
V_S=I^{\otimes b(S)}.
\end{equation}
The differential $d_{V_S}$ is the tensor differential, with sign equal to the sum of the degrees of preceding factors.

For $e\notin S$, the cut $e$ separates two adjacent blocks, say in positions $i,i+1$.
Define
\[
\mu_{S,e}:V_S\longrightarrow V_{S\cup\{e\}},\qquad
 a_1\otimes\cdots\otimes a_{b(S)}
 \longmapsto
 a_1\otimes\cdots\otimes a_i a_{i+1}\otimes\cdots\otimes a_{b(S)}.
\]
This map has degree zero and keeps the order of every surviving block.
It takes values in the augmentation ideal because the augmentation is multiplicative.
The derivation rule for $d_H$ says precisely that it commutes with the tensor differentials.

\begin{lemma}[The interval coefficient squares]
\label{lem:interval-coefficient-squares}
For distinct remaining cuts $e,f$, the two maps to $V_{S\cup\{e,f\}}$ agree:
\[
\mu_{S\cup\{e\},f}\mu_{S,e}
=\mu_{S\cup\{f\},e}\mu_{S,f}.
\]
\end{lemma}

\begin{proof}
If the two cuts have no block in common, each composite multiplies the same two disjoint pairs in their unchanged positions.
Multiplication has degree zero, so interchanging the order of these operations introduces no sign.
If the cuts share a block, three consecutive block states $a,b,c$ are affected.
The two resulting states are $(ab)c$ and $a(bc)$.
They agree by associativity, and all other factors remain unchanged.
These exhaust the possibilities for two distinct cuts in an ordered string.
\end{proof}

Let $U$ be a projected normal chart with one labeled divisor for each cut in $E$.
Using the full residue sheaves on every stratum, define
\begin{equation}
\label{eq:state-residue-carrier}
\mathcal R(V)=\bigoplus_{S\subseteq E}
 i_{S,*}\bigl(\underline{V_S}\otimes\Omega^\bullet_{U_S}(\log D^S)\bigr)[-2|S|].
\end{equation}
The underline denotes the constant sheaf associated to the graded vector space $V_S$.
For homogeneous $v\in V_S$ and a form $\alpha$ on $U_S$, set
\begin{equation}
\label{eq:state-residue-differential}
D_V(v\otimes\alpha)=d_{V_S}v\otimes\alpha
+(-1)^{|v|}v\otimes d\alpha
+\sum_{e\notin S}(-1)^{|v|}\mu_{S,e}(v)\otimes r_{S,e}\alpha.
\end{equation}
The form residue already contains the determinant deletion sign.
The only sign inserted here is the tensor sign from moving it past $v$.
A term has total degree $|v|+|\alpha|+2|S|$.
All three parts of $D_V$ have degree one.

\begin{proposition}[A residue complex with changing ordered states]
\label{prop:state-residue-complex}
The differential~\eqref{eq:state-residue-differential} squares to zero.
Adapted coordinate changes preserving the labeled cut divisors induce strict chain isomorphisms by pullback of forms and the identity on the corresponding $V_S$.
Permuting coordinate names only relabels the same cut diagram and its strata; it does not permute the ordered block states.
These coordinate maps compose strictly.
\end{proposition}

\begin{proof}
The internal and form differentials square to zero and anticommute by the tensor sign.
A collision map commutes with the coefficient differential because $\mu_{S,e}$ is a cochain map.
Its two cross-terms with that differential have signs $(-1)^{|v|}$ and $(-1)^{|v|+1}$, and cancel.
The cross-terms with the form differential cancel by $r_{S,e}d=-dr_{S,e}$.
For two different residues, the tensor signs multiply to $(-1)^{2|v|}=1$.
The coefficient composites agree by Lemma~\ref{lem:interval-coefficient-squares}, while the form composites have opposite signs.
Repeated-edge contributions are absent because their targets have already lost that normal direction.
This proves $D_V^2=0$.

Pullback commutes with the form differential and with each typed residue by Proposition~\ref{prop:strict-residue-descent}.
It acts as the identity on states, so it also commutes with every $d_{V_S}$ and $\mu_{S,e}$.
This proves the chain identity and, by composition of pullbacks, the cocycle.
A relabeling of coordinate slots sends the stratum name to the corresponding set of intrinsic cuts.
It makes no change to that set of blocks or to their order, proving the permutation assertion.
\end{proof}

For three positions and degree-zero states, the four coefficient spaces are
$I^{\otimes3}$, two copies of $I^{\otimes2}$, and $I$.
Applied to $a\otimes b\otimes c\otimes\lambda_1\wedge\lambda_2$, the two residue terms are
\[
(ab)\otimes c\otimes\lambda_2\big|_{D_1},\qquad
-\,a\otimes(bc)\otimes\lambda_1\big|_{D_2}.
\]
Their next residues are $(ab)c$ and $-a(bc)$ on $D_{12}$.
This is the actual merged-state target; associativity cancels the two terms there.
For a noncommutative algebra the order of $a,b,c$ is retained throughout.

These $\mu_{S,e}$ are exactly the degree-zero adjacent multiplication maps of
Proposition~\ref{prop:bar-determinant-comparison}.
If $r$ cuts remain, there are $r+1$ blocks.
The ordinary bar comparison retains the final vertex line in addition to the determinant of those $r$ cuts.
Its determinant therefore has degree $-(r+1)$.
In~\eqref{eq:state-residue-carrier}, the geometric degree instead uses the stratum shift $[-2|S|]$.
Matching the multiplication maps does not identify these two gradings or their entire carriers.

\subsection{The normal embedding for the changing states}
\label{subsec:state-normal-embedding}

On the supplied product projection to $Z$, form
\[
\mathcal A(V)=
\left(\bigoplus_{S\subseteq E}L(E\setminus S)\otimes\underline{V_S}\otimes\Omega_Z^\bullet\right)[-2m].
\]
For $F=E\setminus S$, $r=|F|$, and $q=|v|$, its differential is
\begin{align}
\label{eq:state-normal-differential}
D_{A(V)}(\omega_F\otimes v\otimes\beta)
={}&(-1)^r\omega_F\otimes d_{V_S}v\otimes\beta
+(-1)^{r+q}\omega_F\otimes v\otimes d_Z\beta\notag\\
&+\sum_{e\in F}\iota_e\omega_F\otimes\mu_{S,e}(v)\otimes\beta.
\end{align}
The normal embedding, now with coefficients first in its target, is
\begin{equation}
\label{eq:state-normal-embedding}
j_x^V(\omega_F\otimes v\otimes\beta)
=(-1)^{rq}v\otimes\lambda_F^x\wedge p_S^*\beta.
\end{equation}
The sign moves the $r$ odd determinant generators past $v$.

\begin{proposition}[Comparison for the interval coefficient system]
\label{prop:interval-normal-comparison}
The differential~\eqref{eq:state-normal-differential} squares to zero, and~\eqref{eq:state-normal-embedding}
is an injective degree-zero chain map into $p_*\mathcal R(V)$.
Its image is the direct sum of the corresponding normal-form subspaces over every $S\subseteq E$,
and this image is closed under $D_V$.
When logarithms of the normal coordinate units are chosen, the filtered homotopy is
\begin{equation}
\label{eq:state-logarithmic-homotopy}
K_{xy}^V(\omega_F\otimes v\otimes\beta)
=(-1)^{(r-1)q}v\otimes K_{xy}(\omega_F\otimes\beta),
\end{equation}
where the scalar expression on the right lies on the same stratum $U_S$.
It satisfies $D_VK_{xy}^V+K_{xy}^VD_{A(V)}=j_y^V-j_x^V$.
\end{proposition}

\begin{proof}
The coefficient squares commute, and the coefficient maps are cochain maps.
Opposite determinant signs cancel the two-edge terms in $D_{A(V)}^2$;
the internal signs before and after deletion are $(-1)^r$ and $(-1)^{r-1}$, so those mixed terms cancel as well.
The coefficient differential on $V_S\otimes\Omega_Z$ squares to zero.
Thus $D_{A(V)}^2=0$.

The source degree is $q+|\beta|-r+2m$; the target degree is
$q+|\beta|+r+2|S|$, the same number.
For the coefficient derivative, the sign on the image of the first term of~\eqref{eq:state-normal-differential} is
$(-1)^{r+r(q+1)}=(-1)^{rq}$, as required.
For the base derivative, both paths have sign $(-1)^{rq+r+q}$.
For deletion of the $i$th edge, the signs are
$(-1)^{i-1+(r-1)q}$ and $(-1)^{i-1+rq+q}$; their exponents differ by $2q$.
The coefficient map is $\mu_{S,e}$ on both paths, with the same target.
These calculations prove the chain equation and closure.
Coefficient extraction in the adapted exterior basis, followed by the zero section of the projection,
proves injectivity just as in Lemma~\ref{lem:normal-embedding}.

For the homotopy, substitute the closed parameter forms $\Lambda_e(t)$ of Section~\ref{sec:filtered-homotopies}
for the normal forms in~\eqref{eq:state-normal-embedding}.
Their residues give the same coefficient maps $\mu_{S,e}$.
Extracting a $dt$ and integrating introduces one tensor sign $(-1)^q$ in addition to $(-1)^{rq}$;
this is $(-1)^{(r-1)q}$.
The identity $dI+Id=\operatorname{ev}_1-\operatorname{ev}_0$ and the residue identity $RI+IR=0$
then give the homotopy equation, including the coefficient differential by the same tensor cancellation as in
Lemma~\ref{lem:tensor-residue-homotopies}.
The map stays on $U_S$, so it preserves stratum depth.
\end{proof}

\begin{example}[The normal image and the full stratum summands]
Take $m=1$, $H=\mathbb C[\epsilon]/(\epsilon^2)$ in degree zero, and $Z$ a point.
Then $I=\mathbb C\epsilon$, $V_\varnothing=\mathbb C(\epsilon\otimes\epsilon)$,
and $V_{\{1\}}=\mathbb C\epsilon$.
On the ambient chart the normal image is spanned by
$(\epsilon\otimes\epsilon)\otimes d\log x$, of total degree $1$.
The full ambient summand also contains
$(\epsilon\otimes\epsilon)\otimes1$, of total degree $0$,
which is outside this image.
On the divisor, the normal image is spanned by $\epsilon\otimes1$, of total degree $2$.
Thus both subsets occur in the image, although the ambient summand is not exhausted.
The residue face vanishes because $\epsilon^2=0$; all these normal generators are closed.
\end{example}

An unfiltered homotopy on every unit overlap does not follow for this coefficient diagram.
Take $n=2$, $H=\mathbb C[\epsilon]/(\epsilon^2)$, and $Z=\mathbb C^*$ with coordinate $z$.
Then $V_\varnothing=\mathbb C(\epsilon\otimes\epsilon)$, $V_{\{1\}}=\mathbb C\epsilon$,
and $\mu_{\varnothing,1}=0$.
The coefficient differentials vanish.
On $U=\Delta\times\mathbb C^*$, change the normal coordinate by $y=zx$.
The normal generator $a=e\otimes(\epsilon\otimes\epsilon)\otimes1$ has total degree one and $D_{A(V)}a=0$.
If any degree-minus-one homotopy $G$ existed, $Ga$ would have degree zero in $\mathcal R(V)$.
The collision stratum has minimum degree two, so $Ga=(\epsilon\otimes\epsilon)\otimes f$ on the ambient stratum.
The homotopy equation would imply
\[
df=dz/z.
\]
Restriction to $x=0$ and integration around the unit circle contradict this equation, exactly as in
Example~\ref{ex:period-obstruction}.
Thus no such unfiltered $G$ exists on this overlap, even though the full state-valued residue sheaf has strict coordinate descent.
The average wedge operator $h$ fails here because it would need to reverse the zero multiplication map.

The constructions above give precise carriers for ordered associative states and for a fixed bar complex.
For a chiral product, its input is the sheaf of states away from a collision and its output is supported on the collision diagonal.
It is not the degree-zero multiplication $I\otimes I\to I$ used above.
Applying the displayed formulas to that product requires actual coefficient complexes on every stratum,
a differential-form action, and typed collision maps commuting with the internal differential and the coordinate transitions.
The two coefficient composites at an adjacent double collision must land in the same merged-state complex and satisfy the requisite identity there.
A fixed coefficient comparison additionally requires the stated identification of those coefficient systems.
These are the mathematical conditions for transfer to Definition~\ref{def:bar-differential-complete};
neither determinant signs nor strict descent of logarithmic forms supplies them.
